\chapter{程序如何从源码变成正在运行的进程}

\section{本章要解决的问题}

刚学完 C/C++ 时，你通常知道三件事：

\begin{enumerate}
  \item 写一个 \code{.cpp} 文件。
  \item 用编译器生成可执行文件。
  \item 在终端运行它。
\end{enumerate}

但你可能还不知道：

\begin{itemize}
  \item 头文件到底是不是被“编译进来”。
  \item \code{g++ main.cpp -o app} 背后调用了哪些工具。
  \item \code{main} 为什么通常不是程序真正执行的第一条指令。
  \item 可执行文件里除了机器码还有什么。
  \item 操作系统如何把文件变成一个进程。
  \item 指针地址为什么每次运行可能不同。
  \item 这些机制为什么会影响性能测量。
\end{itemize}

\begin{keyidea}
程序运行不是“源码直接交给 CPU”。源码要经过预处理、编译、优化、汇编、链接、加载和运行时启动，最后 CPU 才开始执行机器指令。性能工程的很多错误判断，都来自看不见这条链路。
\end{keyidea}

\section{从一个最小程序开始}

先看一个最小但真实的 C++ 程序：

\begin{lstlisting}[language=C++]
#include <iostream>

int main()
{
    int x = 1;
    int y = 2;
    std::cout << x + y << '\n';
    return 0;
}
\end{lstlisting}

从人的角度看，它做了三件事：

\begin{center}
\begin{minipage}{0.65\linewidth}
\begin{verbatim}
创建两个整数
计算 x + y
打印结果
\end{verbatim}
\end{minipage}
\end{center}

从系统角度看，它至少涉及：

\begin{center}
\begin{minipage}{0.82\linewidth}
\begin{verbatim}
source file
  -> preprocessor
  -> compiler frontend
  -> optimizer
  -> compiler backend
  -> assembly
  -> assembler
  -> object file
  -> linker
  -> executable ELF
  -> loader
  -> process address space
  -> runtime startup
  -> main
  -> CPU executes instructions
\end{verbatim}
\end{minipage}
\end{center}

本章的任务就是把这条路径铺开。

\section{为什么这条路径必须学}

你最终的目标是高性能计算和 AI CPU 算子优化。为什么第一个基础问题不是 AVX2，而是“程序怎么运行”？

因为很多性能现象发生在源码之外：

\begin{itemize}
  \item 编译器可能删除你以为正在测量的循环。
  \item 函数可能被 inline，二进制里不再有独立函数体。
  \item 动态链接和全局构造可能影响第一次运行时间。
  \item 不同编译参数可能生成完全不同的指令。
  \item 同一个地址表达式可能映射到栈、堆、全局区或 mmap 区域。
  \item 共享库调用、PLT/GOT、lazy binding 可能影响调用路径。
  \item 程序启动成本不是 kernel 成本。
\end{itemize}

如果你不知道这些层，就很容易把系统行为误认为算法行为。

\section{第一层：源文件、头文件和翻译单元}

C/C++ 源文件常见扩展名：

\begin{center}
\begin{verbatim}
.c
.cc
.cpp
.cxx
\end{verbatim}
\end{center}

头文件常见扩展名：

\begin{center}
\begin{verbatim}
.h
.hpp
\end{verbatim}
\end{center}

初学者容易以为头文件会被单独编译，然后和源文件“合并”。通常不是这样。C/C++ 的基本编译单位是 \term{翻译单元}{translation unit}：一个源文件经过预处理后的结果。

如果你写：

\begin{lstlisting}[language=C++]
#include <vector>
#include "kernel.hpp"
\end{lstlisting}

预处理器会把这些头文件的内容展开到当前源文件中。最终编译器看到的是一个很大的文件。

\subsection{为什么头文件模型影响性能}

头文件不是小事。它影响：

\begin{itemize}
  \item inline 函数是否对编译器可见。
  \item 模板代码在哪里实例化。
  \item 宏是否改变源码。
  \item 条件编译是否选择了不同 ISA 路径。
  \item 编译时间和代码膨胀。
\end{itemize}

例如，一个高性能库可能在头文件里写：

\begin{lstlisting}[language=C++]
#if defined(__AVX2__)
// AVX2 implementation
#else
// scalar fallback
#endif
\end{lstlisting}

你以为自己读的是同一份代码，但不同编译参数会选择完全不同的实现。

\section{第二层：预处理器}

\term{预处理器}{preprocessor} 处理 include directive、宏、条件编译和注释删除。

例子：

\begin{lstlisting}[language=C++]
#define SCALE 4

int f(int x)
{
    return x * SCALE;
}
\end{lstlisting}

预处理后大致变成：

\begin{lstlisting}[language=C++]
int f(int x)
{
    return x * 4;
}
\end{lstlisting}

命令：

\begin{lstlisting}[language=bash]
mkdir -p scratch/ch01
cat > scratch/ch01/preprocess.cpp <<'EOF'
#include <iostream>

#define SCALE 4

int f(int x)
{
    return x * SCALE;
}

int main()
{
    std::cout << f(3) << '\n';
}
EOF

g++ -E scratch/ch01/preprocess.cpp -o scratch/ch01/preprocess.i
wc -l scratch/ch01/preprocess.cpp scratch/ch01/preprocess.i
\end{lstlisting}

你通常会看到 \code{preprocess.i} 远比原文件长，因为标准库头文件展开了大量声明和模板。

\begin{warningbox}
不要以为编译器看到的就是你手写的 \code{.cpp}。宏、include 和条件编译会先改写输入。性能分析遇到不符合直觉的行为时，必要时要看预处理结果。
\end{warningbox}

\section{第三层：编译器前端}

预处理结果仍然是 C++ 文本。编译器前端要把它转成内部结构。

简化流程：

\begin{center}
\begin{minipage}{0.78\linewidth}
\begin{verbatim}
characters
  -> tokens
  -> parser
  -> AST
  -> semantic analysis
  -> IR generation
\end{verbatim}
\end{minipage}
\end{center}

\subsection{Token}

代码：

\begin{lstlisting}[language=C++]
int x = a + 3;
\end{lstlisting}

会被词法分析成类似：

\begin{center}
\begin{verbatim}
int
x
=
a
+
3
;
\end{verbatim}
\end{center}

\subsection{AST}

\term{抽象语法树}{Abstract Syntax Tree} 描述语法结构。上面的声明可理解成：

\begin{center}
\begin{minipage}{0.68\linewidth}
\begin{verbatim}
declaration x
  type: int
  initializer:
    binary +
      left: a
      right: 3
\end{verbatim}
\end{minipage}
\end{center}

\subsection{语义分析}

语义分析检查：

\begin{itemize}
  \item 名字是否声明。
  \item 类型是否匹配。
  \item 函数调用参数是否正确。
  \item 重载解析选择哪个函数。
  \item 模板实例化是否合法。
  \item \code{const}、引用和生命周期规则是否满足。
\end{itemize}

这一步决定了程序的 C++ 语义。后续优化必须在这个语义边界内进行。

\section{第四层：IR 和优化器}

编译器通常不会直接从 C++ 生成最终机器码，而是先生成 \term{中间表示}{intermediate representation}，简称 IR。

IR 的作用是提供一个比 C++ 简单、比机器码抽象的优化平台。

例如：

\begin{lstlisting}[language=C++]
extern "C" int add_i32(int a, int b)
{
    return a + b;
}
\end{lstlisting}

LLVM IR 可能类似：

\begin{lstlisting}
define i32 @add_i32(i32 %a, i32 %b) {
entry:
  %sum = add i32 %a, %b
  ret i32 %sum
}
\end{lstlisting}

优化器会在 IR 上做：

\begin{itemize}
  \item 删除无用代码。
  \item 常量折叠。
  \item 函数内联。
  \item 循环优化。
  \item 自动向量化。
  \item 别名分析。
\end{itemize}

\subsection{一个重要例子：无用计算会消失}

如果你写：

\begin{lstlisting}[language=C++]
int sum(int n)
{
    int s = 0;
    for (int i = 0; i < n; ++i) {
        s += i;
    }
    return s;
}

int main()
{
    sum(1000000);
}
\end{lstlisting}

如果返回值没有可观察用途，优化器可能删除整个调用。  
这就是 benchmark 必须防止 dead-code elimination 的原因。

\section{第五层：后端和汇编}

编译器后端把优化后的 IR 降低到目标机器。

后端要解决：

\begin{itemize}
  \item 指令选择：使用哪些 x86-64 指令。
  \item 寄存器分配：值放在哪些寄存器。
  \item 指令调度：指令如何排序。
  \item 调用约定：参数和返回值放哪里。
  \item 汇编输出：生成 \code{.s} 文件。
\end{itemize}

生成汇编：

\begin{lstlisting}[language=bash]
g++ -O3 -S -masm=intel scratch/ch01/add.cpp -o scratch/ch01/add.s
\end{lstlisting}

汇编是机器码的人类可读表示。CPU 最终执行的不是汇编文本，而是汇编器编码后的字节。

\section{第六层：汇编器和目标文件}

\term{汇编器}{assembler} 把汇编文本变成目标文件。

命令：

\begin{lstlisting}[language=bash]
g++ -c -O3 scratch/ch01/add.cpp -o scratch/ch01/add.o
\end{lstlisting}

目标文件中包含：

\begin{itemize}
  \item 机器码。
  \item 符号表。
  \item section。
  \item relocation。
  \item 可选调试信息。
\end{itemize}

它通常还不能单独运行，因为外部函数和最终地址还没处理。

\section{第七层：链接器}

一个真实程序通常由多个目标文件和库组成。

例子：

\begin{lstlisting}[language=C++]
// add.cpp
extern "C" int add_i32(int a, int b)
{
    return a + b;
}
\end{lstlisting}

\begin{lstlisting}[language=C++]
// main.cpp
extern "C" int add_i32(int, int);

int main()
{
    return add_i32(1, 2);
}
\end{lstlisting}

构建：

\begin{lstlisting}[language=bash]
g++ -c add.cpp -o add.o
g++ -c main.cpp -o main.o
g++ add.o main.o -o app
\end{lstlisting}

\term{链接器}{linker} 负责把 \code{main.o} 里对 \code{add_i32} 的引用解析到 \code{add.o} 中的定义。

它还会处理：

\begin{itemize}
  \item 标准库。
  \item 启动代码。
  \item 静态库。
  \item 共享库引用。
  \item relocation。
  \item 符号可见性。
\end{itemize}

\section{第八层：ELF 可执行文件}

Linux x86-64 上的可执行文件通常是 \term{ELF}{Executable and Linkable Format}。

ELF 不是“一段纯机器码”。它是一个带元数据的二进制容器，包含：

\begin{itemize}
  \item ELF header。
  \item program headers。
  \item sections。
  \item symbol table。
  \item relocation information。
  \item dynamic linking information。
  \item debug information。
\end{itemize}

查看：

\begin{lstlisting}[language=bash]
file app
readelf -h app
readelf -l app
readelf -S app
nm -C app
objdump -drwC -Mintel app
\end{lstlisting}

现在不要求你完全读懂每一项。第 6 章会系统展开 ELF。这里你只需要知道：可执行文件是操作系统加载器能理解的一种结构化文件。

\section{第九层：操作系统加载程序}

当你运行：

\begin{lstlisting}[language=bash]
./app
\end{lstlisting}

shell 会请求内核创建新进程。内核和加载器大致做：

\begin{center}
\begin{minipage}{0.82\linewidth}
\begin{verbatim}
read ELF header
create process address space
map executable segments
map shared libraries if needed
prepare stack
prepare argc / argv / envp
start dynamic linker if dynamically linked
transfer control to entry point
\end{verbatim}
\end{minipage}
\end{center}

如果是动态链接程序，动态链接器还会解析共享库和符号，然后运行 C/C++ runtime 初始化。

\begin{keyidea}
\code{main} 是 C/C++ 程序员约定的入口函数，但进程真正开始执行的位置通常是 ELF header 中的 entry point，例如 \code{_start}。
\end{keyidea}

\section{第十层：进程和虚拟地址空间}

\term{进程}{process} 可以理解为：

\begin{center}
\begin{minipage}{0.75\linewidth}
\begin{verbatim}
running program instance
  + virtual address space
  + OS-managed resources
\end{verbatim}
\end{minipage}
\end{center}

进程拥有：

\begin{itemize}
  \item 虚拟地址空间。
  \item 文件描述符。
  \item 当前工作目录。
  \item 信号处理状态。
  \item 一个或多个线程。
  \item 权限和资源限制。
\end{itemize}

简化的虚拟地址空间：

\begin{center}
\begin{minipage}{0.5\linewidth}
\begin{verbatim}
high address
  stack
  ...
  shared libraries / mmap
  ...
  heap
  bss
  data
  rodata
  text
low address
\end{verbatim}
\end{minipage}
\end{center}

这只是模型。真实地址受 ASLR、动态链接、内核配置和运行环境影响。

\section{入口点不是 \code{main}}

查看入口点：

\begin{lstlisting}[language=bash]
readelf -h app | rg "Entry point"
objdump -d -Mintel app | sed -n '/<_start>:/,+40p'
nm -C app | rg " main$|_start"
\end{lstlisting}

你会看到 \code{_start} 一类符号。它最终会调用运行时库，再调用你的 \code{main}。

这对性能测量很重要：

\begin{itemize}
  \item 程序启动成本不是函数成本。
  \item 动态链接和首次解析可能影响第一次运行。
  \item 全局对象构造可能在 \code{main} 前发生。
  \item benchmark 应该把 setup 和 timed region 分开。
\end{itemize}

\section{贯穿实验：拆开一个程序}

\begin{labbox}
创建文件：

\begin{lstlisting}[language=bash]
mkdir -p scratch/ch01
cat > scratch/ch01/add.cpp <<'EOF'
extern "C" int add_i32(int a, int b)
{
    return a + b;
}
EOF
\end{lstlisting}

逐层生成中间产物：

\begin{lstlisting}[language=bash]
g++ -E scratch/ch01/add.cpp -o scratch/ch01/add.i
g++ -S -O0 -masm=intel scratch/ch01/add.cpp -o scratch/ch01/add_O0.s
g++ -S -O3 -masm=intel scratch/ch01/add.cpp -o scratch/ch01/add_O3.s
g++ -c -O3 scratch/ch01/add.cpp -o scratch/ch01/add.o
objdump -drwC -Mintel scratch/ch01/add.o
\end{lstlisting}

观察任务：

\begin{enumerate}
  \item \code{.i}、\code{.s}、\code{.o} 分别是什么。
  \item \code{-O0} 和 \code{-O3} 汇编差异。
  \item \code{objdump} 输出中的机器码字节。
  \item 函数符号名是否保留为 \code{add_i32}。
\end{enumerate}
\end{labbox}

\section{实验：观察地址空间}

\begin{labbox}
\begin{lstlisting}[language=bash]
cat > scratch/ch01/address.cpp <<'EOF'
#include <iostream>
#include <vector>

int global_initialized = 42;
int global_zero;
char const* message = "hello";

int main()
{
    int stack_value = 1;
    auto* heap_value = new int(2);
    std::vector<int> vec(16, 3);

    std::cout << "main               " << reinterpret_cast<void*>(&main) << '\n';
    std::cout << "global_initialized " << &global_initialized << '\n';
    std::cout << "global_zero        " << &global_zero << '\n';
    std::cout << "message variable   " << &message << '\n';
    std::cout << "message literal    " << static_cast<void const*>(message) << '\n';
    std::cout << "stack_value        " << &stack_value << '\n';
    std::cout << "heap_value         " << heap_value << '\n';
    std::cout << "vector data        " << vec.data() << '\n';

    delete heap_value;
}
EOF

g++ -O0 -g scratch/ch01/address.cpp -o scratch/ch01/address
./scratch/ch01/address
./scratch/ch01/address
\end{lstlisting}

报告任务：

\begin{enumerate}
  \item 比较两次运行地址是否一致。
  \item 推测哪些属于 text、data、BSS、rodata、stack、heap。
  \item 解释为什么地址可能每次变化。
  \item 查看 \code{/proc/self/maps} 的方式，并解释它能提供什么信息。
\end{enumerate}
\end{labbox}

\section{反例：看似合理但错误的理解}

\begin{warningbox}
\textbf{误区一：C++ 程序从 \code{main} 的第一行开始。}

不准确。进程入口通常是 runtime 启动代码，\code{main} 是被运行时调用的用户入口。

\textbf{误区二：头文件会被单独编译。}

通常不对。头文件被包含进源文件，形成翻译单元。

\textbf{误区三：目标文件就是可执行文件。}

不对。目标文件还需要链接。

\textbf{误区四：源码里有函数，机器码里一定有这个函数。}

不对。优化后函数可能被 inline、删除、合并或改变可见性。

\textbf{误区五：指针就是物理地址。}

用户态程序看到的是虚拟地址。
\end{warningbox}

\section{作业}

\begin{exercisebox}
\textbf{作业 1：画出完整流程。}

画出从 \code{hello.cpp} 到正在运行进程的完整路径。每一层写明：

\begin{itemize}
  \item 输入是什么。
  \item 输出是什么。
  \item 哪个工具或系统组件负责。
  \item 性能学习为什么要关心。
\end{itemize}

\textbf{作业 2：函数消失实验。}

写三个函数：

\begin{itemize}
  \item 一个允许 inline。
  \item 一个使用 \code{noinline} 阻止 inline。
  \item 一个从不调用。
\end{itemize}

分别用 \code{-O0} 和 \code{-O3} 编译，用 \code{nm} 和 \code{objdump} 观察哪些符号存在，哪些函数消失。

\textbf{作业 3：启动成本和 kernel 成本。}

写一个程序，让 \code{main} 调用某个小函数一次，再调用同一函数一亿次。分别测整个程序时间和循环区域时间。解释为什么启动成本不能代表 kernel 成本。
\end{exercisebox}

\section{验收标准}

完成本章后，你应该能：

\begin{itemize}
  \item 解释预处理、编译、汇编、链接、加载的区别。
  \item 生成并解释 \code{.i}、\code{.s}、\code{.o} 和可执行文件。
  \item 用 \code{readelf}、\code{nm}、\code{objdump} 做基础观察。
  \item 解释为什么 \code{main} 不是真正的进程入口。
  \item 画出简化虚拟地址空间。
  \item 说明为什么性能测量要排除启动、初始化、动态链接和全局构造等因素。
\end{itemize}
